\documentclass[11pt]{article}
\usepackage[margin=2.6cm]{geometry}
\usepackage{amsmath,amssymb,amsthm}
\usepackage{parskip}
\usepackage[T1]{fontenc}
\usepackage{booktabs}

\newcommand{\Hh}{\mathcal{H}}
\newcommand{\C}{\mathcal{C}}
\newcommand{\Rthree}{{}^{(3)}\!R}
\newcommand{\Lie}{\pounds}

\title{The ADM formalism: constraints, multipliers, and what the lapse can and cannot do\\
\large A corrected narrative, with the dictionary to the triangulation program}
\author{}
\date{2026--07--16}

\begin{document}
\maketitle

\section{The story, graded}

The narrative under examination: \emph{begin with a Lorentzian 4-manifold $(M,g)$ and a
foliation by spacelike hypersurfaces $\Sigma_t$; the intrinsic curvature, extrinsic
curvature, lapse, and shift are defined by the foliation; imposing the Einstein equations
forces a local constraint; and} --- the flagged step --- \emph{one can force the constraint
to vanish by choosing the lapse and shift appropriately, for any intrinsic and extrinsic
curvatures, because they act like Lagrange multipliers.}

Everything is correct up to the flagged step, with one refinement. The foliation (together
with $g$ and a labeling $t$ of the leaves) determines the induced metric $h_{ij}$, the
extrinsic curvature $K_{ij}$, and the lapse $N$ (normal proper-time spacing between
leaves: $d\tau = N\,dt$). The \textbf{shift $N^i$ requires one further choice} --- a
\emph{threading}, i.e.\ an identification of points on successive leaves, equivalently a
time-flow vector $t^\mu$ with $t^\mu\partial_\mu t = 1$, decomposed as
\[
t^\mu \;=\; N\,n^\mu \;+\; N^\mu, \qquad N^\mu n_\mu = 0 .
\]
Foliation = the slicing; threading = how the spatial grid rides it.

The flagged step is \textbf{exactly backwards}, in an instructive way.

\section{The correction: the constraints do not contain the lapse}

In vacuum (matter terms shown for completeness) the constraints read
\begin{align}
\Hh_\perp &\;=\; \Rthree + K^2 - K_{ij}K^{ij} \;=\; 16\pi G\,\rho ,\\
\Hh^i &\;=\; D_j\!\left(K^{ij} - h^{ij}K\right) \;=\; 8\pi G\, j^i ,
\end{align}
with $K = h^{ij}K_{ij}$ and $D$ the Levi-Civita connection of $h$. \textbf{The lapse and
shift appear nowhere.} The constraints are conditions purely on the leaf data
$(h_{ij}, K_{ij})$ and the matter data. Given arbitrary $(h,K)$, the constraints are
satisfied or they are not; no choice of $(N, N^i)$ can change that.

A Lagrange multiplier does precisely the opposite of what the flagged step attributes to
it: \emph{a multiplier enforces a constraint on the other variables; it cannot satisfy the
constraint itself, because the constraint does not contain it.} In ordinary mechanics,
\[
L \;=\; p\,\dot q - H(p,q) - \lambda\, C(p,q):
\]
varying $\lambda$ yields $C(p,q) = 0$, a condition on $(p,q)$ in which $\lambda$ is
absent.

\textbf{Warm-up example (electromagnetism).} Gauss's law $\nabla\!\cdot\!\mathbf{E} = \rho$
constrains the initial \emph{field} data; the scalar potential $A_0$ is the multiplier
that enforces it. One would never say ``choose $A_0$ to make
$\nabla\!\cdot\!\mathbf{E}-\rho$ vanish'' --- $A_0$ is not in the equation. The lapse is
gravity's $A_0$.

\section{The story, corrected}

\subsection{Constraints are embedding integrability conditions (Gauss--Codazzi)}

For \emph{any} hypersurface $\Sigma \hookrightarrow (M,g)$, the Gauss and Codazzi
equations relate the intrinsic curvature of $\Sigma$, its second fundamental form $K$,
and the ambient curvature. Tracing the Gauss equation twice and the Codazzi equation once
gives identities
\[
G_{\mu\nu}n^\mu n^\nu \;=\; \tfrac12\!\left(\Rthree + K^2 - K_{ij}K^{ij}\right),
\qquad
G_{\mu\nu}n^\mu e^\nu_i \;=\; D_jK^j{}_i - D_iK .
\]
Imposing the Einstein equation $G_{\mu\nu} = 8\pi G\,T_{\mu\nu}$ turns the left-hand
sides into matter densities, and the constraints are derived. They are nothing exotic:
\textbf{the Gauss--Codazzi integrability conditions for embedding $(\Sigma,h,K)$ into an
Einstein manifold}. Data violating them is not ``a slice with a bad lapse''; it is not a
slice of \emph{any} Einstein spacetime whatsoever.

\subsection{The ten Einstein equations split $4+6$}

The four projections above ($nn$ and $ne_i$) contain no second time derivatives: they are
relations among data on a single leaf --- constraints. The remaining six ($e_ie_j$) are
the genuine evolution equations. Schematically (up to sign conventions):
\begin{align}
\dot h_{ij} &\;=\; 2N\,K_{ij} + D_iN_j + D_jN_i ,\\
\dot K_{ij} &\;=\; -D_iD_jN
+ N\big(\Rthree_{ij} + K\,K_{ij} - 2K_{ik}K^k{}_j\big)
+ \Lie_{\vec N}\,K_{ij} + \text{(matter)} .
\end{align}
\emph{Here} is where the lapse and shift live: they steer the evolution --- how fast the
geometry marches forward at each point, and how the coordinates slide. The first equation
is really the definition of $K$ rearranged: extrinsic curvature is the velocity of the
3-metric measured in proper time.

\subsection{The multiplier statement, made exact}

In Hamiltonian form the ADM action is
\[
S \;=\; \int dt\, d^3x \;\Big[\pi^{ij}\dot h_{ij}
\;-\; N\,\Hh_\perp(h,\pi) \;-\; N^i\,\Hh_i(h,\pi)\Big],
\]
with $N, N^i$ appearing \textbf{linearly and undifferentiated}. The two halves of
``Lagrange multiplier'':
\begin{enumerate}
\item $\delta S/\delta N = 0 \;\Rightarrow\; \Hh_\perp = 0$: varying the multiplier
\emph{yields} the constraint on $(h,\pi)$.
\item Nothing in the action determines $N, N^i$ themselves: they are freely specifiable
functions.
\end{enumerate}
The corrected version of the flagged step is therefore:
\begin{quote}
\emph{For any constraint-satisfying initial data, the lapse and shift may be chosen
arbitrarily, and every choice evolves the data into the same 4-geometry, merely sliced
and coordinatized differently.}
\end{quote}
The freedom lives in the multipliers; the constraint lives in the data.

\subsection{The consistency miracle: constraints propagate}

The scheme would collapse if evolving with some lapse could \emph{break} the constraints
downstream. It cannot: if $\Hh_\perp = \Hh_i = 0$ initially, they vanish for all time,
for any $(N,N^i)$. On the Lagrangian side this is the contracted Bianchi identity
$\nabla_\mu G^{\mu\nu} \equiv 0$; on the Hamiltonian side, the constraints are
\textbf{first-class} --- their Poisson brackets close on constraints (Dirac's
hypersurface-deformation algebra):
\begin{align}
\{\Hh[N], \Hh[M]\} &\;=\; \Hh_i\big[\,h^{ij}(N\partial_j M - M\partial_j N)\,\big],\\
\{\Hh_i[\vec\xi\,], \Hh[N]\} &\;=\; \Hh[\,\xi^i\partial_i N\,],\qquad
\{\Hh_i, \Hh_j\} \;\sim\; \text{(spatial diffeomorphism algebra)} .
\end{align}
Note the $h^{ij}$ in the first bracket: a \emph{field-dependent structure function}, not
a structure constant --- this is not a Lie algebra. (This is the object the
``deformation-algebra test'' of the lapse program aspires to see, and the structure
Dittrich showed discretization generically breaks.)

\subsection{The payoff: general relativity as an initial-value problem}

Give $(\Sigma, h, K)$ satisfying the four constraints; choose \emph{any} lapse and shift;
evolve. Choquet-Bruhat (existence) and Choquet-Bruhat--Geroch (uniqueness) give the
unique maximal globally hyperbolic Einstein development. Degrees-of-freedom audit: $12$
phase-space functions per point $(h_{ij}, \pi^{ij})$, minus $4$ constraints, minus $4$
gauge orbits they generate, $= 4$: the two polarizations of gravitational waves.

Solving the constraints for admissible initial data is its own elliptic theory (the
York--Lichnerowicz conformal method) --- a likely source of the half-memory behind the
flagged step: one \emph{does} solve elliptic equations to satisfy the constraints, but
the unknowns are a conformal factor and a vector potential inside $(h,K)$, never the
lapse. The lapse's elliptic equation (e.g.\ maximal slicing's $D^2N = N\,K_{ij}K^{ij}$)
arises only \emph{after} the constraints, as an optional gauge condition.

\section{Dictionary to the triangulation program}

\begin{center}
\begin{tabular}{@{}lll@{}}
\toprule
ADM object & program object & the corrected role \\
\midrule
data $(h,K)$ & the triangulation ensemble & constraints condition \emph{it} \\
constraint $\Hh_\perp$ & patch constraint $\C_p$ (tier c) & contains geometry only, never $N_p$ \\
multiplier $N$ & lapse field $N_p$; kernel freedom at fixed $\pi$ & enforces; is itself undetermined \\
constraint propagation & automatic: each local $h_x$ annihilates $\sqrt{\pi}$ & trivially abelian algebra \\
Dirac algebra & (absent; the target) & field-dependent structure functions \\
\bottomrule
\end{tabular}
\end{center}

The reason general relativity must \emph{check} constraint propagation while the sampler
gets consistency for free: the sampler's $H[N] = \sum_x N_x h_x$ has each local generator
$h_x$ \emph{individually} annihilating the state --- a trivially closed, abelian
constraint structure --- whereas gravity's constraints close only jointly, with
field-dependent structure functions. That gap --- trivial algebra versus Dirac algebra
--- is, in one sentence, the distance between ``a sampler with a lapse'' and
``gravity.''

\end{document}
